\section{控制算法设计}

\subsection{预设时间扰动观测器设计}

对等效关节空间模型 \eqref{eq:joint_model}，先设计扰动观测器来生成后续动力学补偿所需的 \(\hat d(t)\)。将关节速度记为
\[
\chi=\dot q_m .
\]
由 \eqref{eq:joint_model} 可得
\begin{equation}
\dot\chi=u_o+\Delta_a,
\qquad
u_o=M_e^{-1}(q_m)\bigl[\tau-C_e(q_m,\dot q_m)\dot q_m\bigr],
\qquad
\Delta_a=M_e^{-1}(q_m)d(t).
\label{eq:observer_channel}
\end{equation}
其中 \(\Delta_a\) 为匹配到加速度通道的集总扰动。该形式与预设时间扰动观测器的标准受扰单积分器模型一致\cite{jiangPrescribedtimeDisturbanceObserver2024}。若能够在预设时间 \(T_o\) 内得到 \(\hat\Delta_a\)，则力矩扰动估计可由
\begin{equation}
\hat d(t)=M_e(q_m)\hat\Delta_a(t)
\label{eq:dhat_from_observer}
\end{equation}
给出。

设 \(z_1,z_2\in\mathbb R^n\) 分别为 \(\chi\) 和 \(\Delta_a\) 的估计值，取 \(\hat\Delta_a=z_2\)。定义可测复合误差
\[
\varepsilon_1=\chi-z_1-\bar\xi(t),
\]
其中 \(\bar\xi(t)\in\mathbb R^n\) 为调节函数，满足有界光滑、单调衰减且 \(t\ge T_o\) 时 \(\bar\xi(t)=0\)。例如可取
\[
\bar\xi(t)=
\begin{cases}
\bar\xi_0(T_o-t)^2, & 0\le t<T_o,\\
0, & t\ge T_o.
\end{cases}
\]
记 \(\operatorname{sig}^{\alpha}(x)=\operatorname{sgn}(x)|x|^\alpha\)，并按元素作用于向量。参考文献\cite{jiangPrescribedtimeDisturbanceObserver2024}，构造如下预设时间扰动观测器：
\begin{equation}
\left\{
\begin{aligned}
\dot z_1
&=z_2+\frac{\pi}{\eta T_c}
\phi_1\left(\frac{\varepsilon_1}{\sigma}\right)
-\dot{\bar\xi}(t)+\bar\xi(t)+u_o,\\
\dot z_2
&=\frac{\pi}{\sigma\eta T_c}
\phi_2\left(\frac{\varepsilon_1}{\sigma}\right)
-\dot{\bar\xi}(t),
\end{aligned}
\right.
\label{eq:ptdo}
\end{equation}
其中 \(T_c\le T_o\)、\(\eta\in(0,1)\)、\(\sigma>0\)，且
\[
\begin{aligned}
\phi_1(x)&=\operatorname{sig}^{1-\eta/2}(x)+\operatorname{sig}^{1+\eta/2}(x),\\
\phi_2(x)&=\operatorname{sig}^{2-\eta}(x)+\operatorname{sig}^{2+\eta}(x)+\sigma\operatorname{sgn}(x).
\end{aligned}
\]
在 \(\Delta_a\) 与 \(\dot\Delta_a\) 有界、连续时间实现且观测器满足相应理论条件时，观测误差可在设计时间 \(T_o\) 内收敛到零。该结论仅适用于连续时间理想观测器；在采样、量化、测量噪声和执行器限幅存在时，观测误差通常只能收敛到与采样步长、数值积分误差和未建模动态有关的小邻域。参数 \(T_o\) 和 \(T_c\) 不宜选得过小，否则可能带来较大的观测峰值和控制输入峰值。

\subsection{动力学补偿与误差通道规范化}

在获得 \(\hat d(t)\) 后，采用如下动力学补偿结构
\begin{equation}
\tau=M_e(q_m)v+C_e(q_m,\dot q_m)\dot q_m-\hat d(t),
\label{eq:computed_torque}
\end{equation}
其中 \(v\) 为辅助加速度输入。将式 \eqref{eq:computed_torque} 代入式 \eqref{eq:joint_model} 可得
\[
M_e(q_m)(\ddot q_m-v)=d(t)-\hat d(t).
\]
当观测器已收敛并满足 \(\hat d(t)=d(t)\) 时，且 \(M_e(q_m)\) 正定可逆，因此
\(
\ddot q_m=v.
\)

令
\(
v=\ddot q_d+\nu ,
\)
并定义
\(
e=q_m-q_d\),
\(z=[e, \dot e]^T.
\)
则
\(
\ddot e=\ddot q_m-\ddot q_d=\nu,
\)
误差系统可写成二阶积分链形式
\begin{equation}
\dot z=Az+B\nu ,
\label{eq:linear_error}
\end{equation}
其中
\[
A=\begin{bmatrix}0&I_n\\0&0\end{bmatrix},\qquad
B=\begin{bmatrix}0\\I_n\end{bmatrix}.
\]
矩阵对 \((A,B)\) 可控，因此 PDF 增益的设计对象由原非线性动力学转移为线性可控误差通道。这个步骤是本文方法成立的关键：PDF 理论只需作用在 \eqref{eq:linear_error} 上，而不直接处理自由漂浮机械臂的非线性耦合项。

\subsection{PDF 线性误差系统设计}

文献\cite{zhouFixedtimeStabilizationLinear2022}考虑的基本对象是可控线性时滞系统，其固定时间镇定依赖如下周期延迟结构：
\begin{equation}
\dot x(t)=F(t)x(t)+G(t)x(t-h),
\label{eq:periodic_delay_general}
\end{equation}
其中 \(F(t),G(t)\) 为 \(2h\) 周期矩阵，且 \(G(t)=0\) 在每个周期的前半段 \([0,h]\) 成立。由分步法可知，系统在一个周期后的状态满足
\(
x(2(k+1)h)=\Delta(h)x(2kh),
\)
其中单周期状态转移矩阵（又称单值矩阵）为
\begin{equation}
\Delta(h)=\Phi(2h,0)+\int_h^{2h}\Phi(2h,s)G_0(s)\Phi(s-h,0)\,ds.
\label{eq:period_map}
\end{equation}
若 \(\Delta(h)\) 为幂零矩阵，即存在最小整数 \(\nu_0\ge1\) 使 \(\Delta^{\nu_0}(h)=0\)，则 \(x(t)=0,\ t\ge 2\nu_0h\)。特别地，若通过增益设计使 \(\Delta(h)=0\)，则系统可在一个完整周期后归零\cite{zhouFixedtimeStabilizationLinear2022}。

为将该思想用于式 \eqref{eq:linear_error}，先取常值反馈 \(K_0\) 使
\(
A_c=A+BK_0
\)
成为期望的基准闭环矩阵，再引入周期延迟反馈
\begin{equation}
\nu(t)=K_0z(t)-K_c(t)z(t-h).
\label{eq:pdf_error_input}
\end{equation}
此时误差闭环为
\begin{equation}
\dot z(t)=A_cz(t)-BK_c(t)z(t-h).
\label{eq:closed_delay_error}
\end{equation}
令 \(R_h(t)\) 为满足
\[
R_h(t+2h)=R_h(t)
\]
的 \(2h\) 周期光滑函数，并满足
\[
R_h(t)=0,\quad t\in[0,h],\qquad
R_h(t)\ge0,\quad t\in[h,2h],
\]
且在 \([h,2h]\) 的某个具有非零测度的子区间上严格为正。本文采用标量函数情形。根据文献\cite{zhouFixedtimeStabilizationLinear2022}的构造，可定义
\begin{equation}
W_c(A_c,h)=\int_h^{2h}e^{-A_cs}BR_h(s)B^\top e^{-A_c^\top s}\,ds,
\label{eq:wc_matrix}
\end{equation}
并取
\begin{equation}
K_c(t)=R_h(t)B^\top e^{-A_c^\top t}W_c^{-1}(A_c,h)e^{A_c(h-t)}.
\label{eq:kc_pdf}
\end{equation}
由于 \((A_c,B)\) 与 \((A,B)\) 同可控，且 \(R_h(t)\) 在后半周期的一个非零测度子区间上严格为正，连续时间条件下 \(W_c(A_c,h)\) 为正定并可逆。将 \eqref{eq:kc_pdf} 代入 \eqref{eq:period_map} 可使闭环的单周期状态转移矩阵满足 \(\Delta(h)=0\)，从而得到单周期零化性质。该结论依赖精确的 Gramian 计算和连续时间延迟实现；数值积分、矩阵病态和采样误差会使实际闭环仅近似满足零化条件。式 \eqref{eq:kc_pdf} 也说明 PDF 增益可通过离线积分计算，在线控制只需读取当前误差和历史误差。

结合式 \eqref{eq:computed_torque} 和式 \eqref{eq:pdf_error_input}，PDF 关节力矩控制律可写为
\begin{equation}
\tau
=
M_e(q_m)
\left[
\ddot q_d
+
K_0z(t)
-
K_c(t)z(t-h)
\right]
+
C_e(q_m,\dot q_m)\dot q_m
-
\hat d(t).
\label{eq:final_controller}
\end{equation}

\subsection{固定时间抗饱和补偿器}

实际应用中，关节执行器存在力矩饱和限制，即
\begin{equation}
\tau_i = \operatorname{sat}(\tau_{ci}) =
\begin{cases}
\tau_{i,\max}\,\operatorname{sgn}(\tau_{ci}), & |\tau_{ci}|>\tau_{i,\max},\\
\tau_{ci}, & |\tau_{ci}|\le\tau_{i,\max},
\end{cases}
\label{eq:saturation_model}
\end{equation}
其中 \(\tau_{ci}\) 为第 \(i\) 个关节的计算控制力矩，\(\tau_{i,\max}>0\) 为力矩幅值上限。饱和将导致实际施加力矩与计算力矩之间产生偏差 \(\Delta\tau_i=\tau_i-\tau_{ci}\)。映射到加速度通道的归一化饱和误差为
\begin{equation}
\Delta_u(t)=M_e^{-1}(q_m)\Delta\tau(t).
\label{eq:norm_saturation_error}
\end{equation}

为抑制输入饱和对跟踪性能的影响，参考文献\cite{dou2025fractional}，引入如下固定时间抗饱和补偿器：
\begin{equation}
\dot\eta = -L_\eta\frac{k_{\eta1}}{\rho_1}
\tanh\bigl(|\eta|^{1-\rho_1}\bigr)
\cosh\bigl(|\eta|^{2}\bigr)
\operatorname{sgn}(\eta)
- L_\eta k_{\eta2}
\operatorname{sig}^{\rho_2}(\eta)
+ \Delta_u,
\label{eq:as_dynamics}
\end{equation}
其中 \(\eta\in\mathbb R^n\) 为辅助变量，\(0<\rho_1,\rho_2<1\)，\(k_{\eta1},k_{\eta2}>0\) 为增益系数，\(L_\eta>0\) 为自适应增益。\(\operatorname{sig}^{\alpha}(x)=\operatorname{sgn}(x)|x|^{\alpha}\) 按元素作用于向量。该补偿器的核心性质由以下定理给出：

\textbf{定理1：} 当 \(\Delta_u=0\) 时，辅助变量 \(\eta\) 将在固定时间内收敛至零，其上界为
\begin{equation}
T_{\eta,\max} < 
\frac{\rho_1 I}{k_{\eta1}\tanh(1)} + 
\frac{1}{k_{\eta2}(1-\rho_2)},
\end{equation}
其中 \(I=\int_1^{\infty}\frac{dz}{\cosh(z^2)}<\infty\) 为收敛常数。

\textbf{证明：} 令 \(z=|\eta|\)，由式 \eqref{eq:as_dynamics} 可得
\[
\dot z = -L_\eta\frac{k_{\eta1}}{\rho_1}
\tanh(z^{1-\rho_1})
\cosh(z^2)
- L_\eta k_{\eta2} z^{\rho_2}.
\]
当 \(z\ge 1\) 时，\(\tanh(z^{1-\rho_1})\ge\tanh(1)\) 且 \(\cosh(z^2)\ge e^{z^2}/2\)，因此
\[
T_3 = \int_1^{z_0}\frac{dz}{|\dot z|}
\le \frac{\rho_1}{L_\eta k_{\eta1}\tanh(1)}
\int_1^{\infty}\frac{dz}{\cosh(z^2)}
= \frac{\rho_1 I}{L_\eta k_{\eta1}\tanh(1)}.
\]
当 \(z\le 1\) 时，\(\dot z\le -L_\eta k_{\eta2}z^{\rho_2}\)，解得
\(
T_4 < 1/(L_\eta k_{\eta2}(1-\rho_2))
\)。
因此总收敛时间上界为 \(T_{\eta,\max}=T_3+T_4\)，与初始值 \(z_0\) 无关。\hfill\(\square\)

将抗饱和补偿器与 PDF 控制器结合，得到如下修正控制律：
\begin{equation}
\begin{aligned}
\tau_c &=
M_e(q_m)\bigl[
\ddot q_d + K_0z(t) - K_c(t)z(t-h) + \eta
\bigr] \\
&\qquad + C_e(q_m,\dot q_m)\dot q_m - \hat d(t).
\label{eq:final_controller_as}
\end{aligned}
\end{equation}
实际施加力矩为 \(\tau=\operatorname{sat}(\tau_c)\)。与原始控制律 \eqref{eq:final_controller} 相比，式 \eqref{eq:final_controller_as} 在辅助加速度中增加了 \(\eta\) 项。当 \(|M_e^{-1}\tau_c|>\tau_{\max}\) 时，饱和误差 \(\Delta_u\) 驱动 \(\eta\) 沿式 \eqref{eq:as_dynamics} 演化；一旦饱和解除，\(\eta\) 在固定时间 \(T_{\eta,\max}\) 内衰减至零，控制律退化为原始的 PDF 形式。

需要注意，若观测器与 PDF 控制器同时从 \(t=0\) 启动，则 \(0\le t<T_o\) 内的观测误差会作为残余扰动进入误差系统，严格 \(2h\) 零化结论不再直接成立。为保持理论闭环清晰，可采用两阶段时序：先在 \(T_o\) 内完成扰动观测，随后以局部时间 \(\vartheta=t-T_o\) 启动 PDF 增益 \(K_c(\vartheta)\)。此时理想连续时间条件下的总预设收敛时间上界为
\[
T_{\mathrm{tot}}=T_o+2h .
\]
在非理想实现中，该时间应理解为进入理论实际有界性分析的起始时间尺度，而非严格的数值归零时刻。

% \subsection{实现说明与鲁棒项}

% 在数值实现中，本文默认按照式 \eqref{eq:wc_matrix} 和式 \eqref{eq:kc_pdf} 离线计算 \(W_c(A_c,h)\) 与 \(K_c(t)\)，并在每个采样时刻按 \(2h\) 周期读取局部时间对应的 PDF 增益。为便于调参和消融对比，程序中还保留工程化平滑延迟项
% \[
% K_0=[-K_p\ -K_d],\qquad
% -K_c(t)=\rho R_h(t)[-K_{ph}\ -K_{dh}],
% \]
% 其中 \(K_p,K_d,K_{ph},K_{dh}\) 为正定对角矩阵，\(\rho>0\) 为延迟反馈幅值系数。该工程化形式仅保留 PDF 的等待--作用周期结构，并不等价于严格的单周期状态转移矩阵零化增益；论文中的理论固定时间结论仍以式 \eqref{eq:kc_pdf} 的 Gramian 型 PDF 增益为前提。

% 若存在观测误差或未建模高频扰动，也可进一步引入饱和型鲁棒补偿
% \[
% \tau_r=-k_d\,\operatorname{sat}(s/\varphi),
% \qquad
% s=\dot e+\Lambda e,
% \]
% 并将 \(\tau_r\) 叠加到 \eqref{eq:final_controller}。其中 \(k_d>0\)、\(\varphi>0\)、\(\Lambda=\Lambda^\top>0\)。需要指出，加入非线性鲁棒补偿后，闭环不再是文献\cite{zhouFixedtimeStabilizationLinear2022}中的纯线性延迟系统，严格固定时间精确归零结论需要单独证明；工程上通常得到对扰动鲁棒的实际固定时间有界结果。
